
\appendix


\begin{table*}[htb]
% \scriptsize
\small
\centering
\caption{Definitions of fine-grained deception labels and their explanations.}
\renewcommand{\arraystretch}{1}  % 行间距
\setlength{\tabcolsep}{6pt}         % 列间距
\rowcolors{2}{gray!10}{white}       % 斑马纹行
\begin{tabular}{>{\bfseries}l m{12.5cm}}
\toprule
Fine-grained Deception Label & \textbf{Explanation} \\
\midrule
Real News & The image and text are consistent and factually aligned, representing authentic news. \\
Image Fabrication & The image is likely manipulated or artificially generated. This includes advanced deepfake techniques, simple edits like splicing or Photoshopping, and fabricated visuals imitating official websites or social media posts. \\
Entity Inconsistency & The key entities (e.g., people, places, organizations) mentioned in the text do not match those depicted in the image, indicating a misalignment between modalities. \\
Event Inconsistency & Although the image and text refer to the same entities, they describe different events. The image may be genuine, but the text misrepresents or over-interprets its context. \\
Time or Space Inconsistency & The image or video is real but misrepresented in time or location---for example, showing past or distant events as if they were current or local. \\
Ineffective Visual Information & The image provides no evidential support for the accompanying text. It typically consists of text-only visuals or irrelevant content, lacking concrete proof such as on-site or event-related imagery. \\
\bottomrule
\end{tabular}

\label{tab:fine_label}
\end{table*}


\section{Fine-grained Deception Labels}
\label{appendix:deception_label}
In addition to the binary labels \textsc{Real} and \textsc{Fake}, the AMG dataset further annotates each fake news sample with a fine-grained deception label. Specifically, the fine-grained labels include \textit{Image Fabrication}, \textit{Entity Inconsistency}, \textit{Event Inconsistency}, and \textit{Time or Space Inconsistency}. The extended deception labels $\tilde{y}_i$ and their corresponding explanations are presented in Table~\ref{tab:fine_label}.

\section{Human Evaluation of LLM Explanations}
\label{appendix:human_verify}
We conduct a human evaluation of LLM-generated explanations. Three annotators independently assess 50 randomly sampled instances, given the multimodal input, its label, and the explanation generated by GPT-4o. They rate (0/1/2): (i) \textbf{Label--Explanation Consistency} (\textsc{Consis.}), whether the explanation supports the label; (ii) \textbf{Multimodal Faithfulness} (\textsc{MM-Faith.}), whether it is grounded in the image and text without hallucinated entities or unsupported claims; and (iii) \textbf{Evidence Usefulness} (\textsc{Useful.}), whether it provides specific, actionable clues rather than generic statements. We report mean scores and inter-annotator agreement (Fleiss' $\kappa$) with majority-vote aggregation. We observe substantial agreement ($\kappa{=}0.58$), and the averaged scores are \textsc{Consis.}=$\mathbf{1.94}$, \textsc{MM-Faith.}=$\mathbf{1.75}$, and \textsc{Useful.}=$\mathbf{1.89}$. These results suggest the explanations are generally reliable and useful for retrieval.


\section{Prompt Templates}
\label{sec:appendix_prompting}
For the final prediction, we adopt two prompting strategies corresponding to the two modes. \textbf{Prompt} \textsc{II} is used in the non-retrieval setting, i.e., for baselines without retrieval or when the dynamic trigger deems external evidence unnecessary. When retrieval is activated, we use \textbf{Prompt} \textsc{III}, which prepends the retrieved positive ($\mathbf{x}^{+}$) and negative ($\mathbf{x}^{-}$) instances as few-shot in-context examples to guide a more informed and robust classification.




\section{Details for Token-Level Support Calculation}
\label{appendix:Lexicons}
The token-level support score, $\tau_{\text{tok}}$, quantifies the semantic alignment of the top-$K$ candidate tokens with the model's predicted label. This is determined via a hierarchical matching process.

% \vspace{1mm}
\noindent\textbf{Step 1: Lexical Matching.} We compare each candidate token $t_i$ against  two predefined label-related lexicons:

\begin{itemize}[leftmargin=4mm]
    \item $\mathcal{L}_{\text{real}}$: \texttt{\{real, genuine, authentic, true, legitimate, realistic, legit, fact, accurate, related, likely, consistent, plausible\}}
    \item $\mathcal{L}_{\text{fake}}$: \texttt{\{fake, missing, false, fabric,  fict, un, mis, fraud, unrelated, fictional, inconsistent\}}
\end{itemize}


If a candidate token $t_i$ exactly matches any entry in either lexicon, it is directly assigned the corresponding label (\textsc{Real} or \textsc{Fake}).

\vspace{1mm}
\noindent\textbf{Step 2: Semantic Similarity Voting.} If $t_i$ is not found in either lexicon, we assess its alignment via semantic similarity. The token is inserted into a prompt template to form a query sentence: ``\texttt{This post is $t_i$.}'' and embedded using the \texttt{all-MiniLM-L6-v2} sentence encoder. It is then compared with reference sentence embeddings constructed from each lexicon:

\begin{itemize}[leftmargin=3mm]
    \item Real references: ``\textit{The post is \{word\} and factually correct.}'' for each word in $\mathcal{L}_{\text{real}}$
    \item Fake references: ``\textit{The post is \{word\} and contains misinformation.}'' for each word in $\mathcal{L}_{\text{fake}}$
\end{itemize}

We compute the mean cosine similarity between the query embedding and each reference group:

\begin{align}
\text{sim}_{\text{real}}(t_i) &= \frac{1}{|\mathcal{L}_{\text{real}}|} \sum_{w \in \mathcal{L}_{\text{real}}} \cos(\mathbf{e}_{\text{query}}, \mathbf{e}_{\text{real}, w}) \\
\text{sim}_{\text{fake}}(t_i) &= \frac{1}{|\mathcal{L}_{\text{fake}}|} \sum_{w \in \mathcal{L}_{\text{fake}}} \cos(\mathbf{e}_{\text{query}}, \mathbf{e}_{\text{fake}, w})
\end{align}
where $\mathbf{e}_{\text{query}}$ denotes the embedding of the query, $\mathbf{e}_{\text{real}, w}$ and $\mathbf{e}_{\text{fake}, w}$ denote the embeddings of the reference token $w$ in the real and fake groups, respectively, and $\cos(\cdot,\cdot)$ denotes cosine similarity.

The token is assigned the label with the higher similarity. This enables generalization to unseen but semantically aligned tokens.

\noindent\textbf{Step 3: Score Computation.} The final token-level support score is defined as the proportion of top-$K$ tokens whose assigned label matches the model's predicted label $y$:

\begin{equation}
\tau_{\text{tok}} = \frac{N_{\text{sup}}}{K},
\end{equation}
where $N_{\text{sup}}$ is the number of supporting tokens that agree with $y$ in either the lexical or semantic stage.



\begin{tcolorbox}[title=Prompt \textsc{II}: Fake News Detection without Retrieval, boxrule=0pt, left=1mm, right=1mm, top=1mm, bottom=1mm, fontupper=\small]
You are a knowledgeable and analytical fact-checking assistant. Your task is to determine whether a social text-image pair is fake.\\

Your response should be either \texttt{The pair is fake because \{explanation of your reasoning\}.} if the text and image present false, misleading, or manipulated content, or \texttt{The pair is real because \{explanation of your reasoning\}.} if the text and image are consistent and factually aligned.\\

Your explanation must be concise and clear, highlighting linguistic, visual, or contextual cues that support your conclusion.

\vspace{1mm}
\textbf{USER}: the image \textit{\textless image\textgreater} and the text \textit{\textless content \textgreater}.

\textbf{ASSISTANT}:
\end{tcolorbox}


\begin{tcolorbox}[title=Prompt \textsc{III}: Fake News Detection with Retrieval, boxrule=0pt, left=1mm, right=1mm, top=1mm, bottom=1mm, fontupper=\small]
You are a knowledgeable and analytical fact-checking assistant. Your task is to determine whether a social text-image pair is fake.\\

Your response should be either \texttt{The pair is fake because \{explanation of your reasoning\}.} if the text and image present false, misleading, or manipulated content, or \texttt{The pair is real because \{explanation of your reasoning\}.} if the text and image are consistent and factually aligned.\\

Your explanation should be concise and clear, highlighting any linguistic, visual, or contextual cues that support your conclusion.

\vspace{1mm}
\textbf{Refer to these examples:}
% \begin{itemize}

\textbf{USER}: the first image \textit{\textless positive\_image\textgreater} and the text \textit{\textless positive\_content\textgreater}.

\textbf{ASSISTANT}: The pair is real because \{positive\_explanation\}

\textbf{USER}: the second image \textit{\textless negative\_image\textgreater} and the text \textit{\textless negative\_content\textgreater}.
    
\textbf{ASSISTANT}: The pair is fake because \{negative\_explanation\}
% \end{itemize}

\vspace{1mm}
\textbf{Now determine the following:}

\textbf{USER}: the third image \textit{[image]} and the text \textit{[content]}.

\textbf{ASSISTANT}:
\end{tcolorbox}

